\subsection{Significado del parámetro \texttt{-np} y sus variantes}

El parámetro \texttt{-np}, cuyo acrónimo proviene de la expresión inglesa \textit{number of processes} (número de procesos), constituye el argumento más crítico e indispensable al momento de invocar el gestor de ejecución \texttt{mpirun} o \texttt{mpiexec}. Su propósito primordial consiste en definir de manera explícita el grado de paralelismo inicial, especificando la cantidad exacta de instancias idénticas del binario que el entorno de ejecución distribuido debe inicializar de forma simultánea \cite{gropp_using_mpi_1999}. Este valor numérico altera directamente la arquitectura lógica de la aplicación, ya que determina el tamaño absoluto o cardinalidad del comunicador global primario, denotado formalmente dentro del estándar como \texttt{MPI\_COMM\_WORLD}\footnote{En MPI, la cardinalidad de un comunicador equivale a su tamaño, es decir, al número total de procesos que lo integran. En el caso de \texttt{MPI\_COMM\_WORLD}, esa cardinalidad coincide con el valor proporcionado a \texttt{-np} o \texttt{-n}; en consecuencia, cada proceso obtiene un rango único en el intervalo $[0, n-1]$ y puede consultar dicho tamaño mediante la rutina \texttt{MPI\_Comm\_size}. Esta relación entre cantidad de procesos, rangos y comunicador global es una de las bases del modelo SPMD en MPI \cite{mpi_forum_2025,pacheco_2011}.} \cite{mpi_forum_2025}.

En el aspecto operativo, la utilización de este modificador exige una asignación posicional estricta, colocándose inmediatamente después del comando de lanzamiento y precediendo a la ruta del archivo ejecutable. Por ejemplo, al ejecutar la instrucción \texttt{mpirun -np 16 ./simulacion}, la biblioteca subyacente coordina la creación de dieciséis procesos independientes, asignándole a cada uno de ellos un identificador entero único conocido como rango (\textit{rank}), el cual oscila en el intervalo cerrado $[0, n-1]$, donde $n$ representa el argumento entero provisto \cite{pacheco_2011}. Dependiendo de la distribución específica de MPI configurada en el clúster, este parámetro cuenta con variantes sintácticas equivalentes aceptadas por la industria, tales como \texttt{-n} o \texttt{--c}, las cuales preservan intacta la misma semántica funcional de aprovisionamiento de tareas \cite{open_mpi_2026}. No obstante, para mantener la máxima portabilidad y el apego estricto al estándar, se recomienda priorizar el uso de \texttt{-n}, ya que la documentación oficial de MPI explicita esta forma sobre la variante histórica \texttt{-np} \cite{gropp_using_mpi_1999}.

La selección del valor entero que acompaña a \texttt{-np} o \texttt{-n} demanda una planificación rigurosa basada en la topología del hardware y en el modelo de descomposición de datos del algoritmo paralelo. Un error metodológico ordinario consiste en fijar de forma arbitraria un número de procesos que no guarde coherencia con los requerimientos geométricos de la aplicación, como ocurre en algoritmos de reducción en árbol o mallas bidimensionales que exigen estrictamente potencias de dos o cuadrados perfectos para evitar desbalances de carga de trabajo \cite{grama_intro_parallel_2003}. Asimismo, si el entero provisto supera la cantidad de ranuras de procesamiento físicas especificadas en el archivo de configuración de nodos (\textit{hostfile}) sin haber habilitado explícitamente banderas de tolerancia a la sobreasignación\footnote{Estas banderas permiten solicitar al gestor de ejecución que acepte más procesos que ranuras físicas disponibles en los nodos asignados. En Open MPI, la opción habitual es \texttt{--oversubscribe}, la cual autoriza la sobreasignación y evita que el lanzamiento falle de inmediato por falta de ranuras. Su uso es válido en escenarios controlados, pero puede degradar el rendimiento porque varios procesos comparten el mismo núcleo o compiten por tiempo de CPU \cite{open_mpi_2026}.}, el gestor de procesos abortará la inicialización de inmediato, lo que en la práctica evidencia que la solicitud excede la capacidad asignada para ese lanzamiento y fuerza al usuario a decidir entre ajustar \texttt{-np} o aceptar la sobreasignación bajo control explícito \cite{open_mpi_2026}.

Como se mencionó anteriormente, además de la planificación del número de procesos, el estándar MPI ofrece mecanismos adicionales que otorgan flexibilidad al momento del lanzamiento. Es importante notar que, en circunstancias ordinarias, los modificadores \texttt{-np} y \texttt{-n} operan bajo una semántica de \comillas{todo o nada}: el gestor de procesos intenta asignar exactamente el número de procesos solicitados y, si no es posible encontrar suficientes ranuras disponibles en el clúster (incluso después de habilitar \texttt{--oversubscribe}), la ejecución falla de inmediato con un error de recursos \cite{open_mpi_2026}. 

En aplicaciones heterogéneas del tipo MPMD (\textit{Multiple Program, Multiple Data}), el parámetro puede utilizarse múltiples veces en una misma ejecución separando los bloques con el símbolo de dos puntos (\texttt{:}), lo que permite indicar una cantidad de procesos específica para cada programa; por ejemplo, \texttt{mpiexec -n 5 programa\_A : -n 10 programa\_B} \cite{mpi_forum_2025}. Sin embargo, cada bloque sigue siendo una solicitud rígida: si las ranuras no están disponibles, el lanzamiento aborta. De manera complementaria, cuando existe incertidumbre acerca de la disponibilidad de recursos en tiempo de ejecución o cuando se desea que el gestor aproveche dinámicamente los nodos disponibles sin fallar si no se alcanza el número exacto solicitado, el modificador \texttt{-soft} ofrece una semántica alternativa y más permisiva\footnote{La sintaxis del modificador \texttt{-soft} varía según la implementación. En Open MPI, se utiliza como \texttt{-npernode <min>:<max>} o directamente \texttt{-soft <rango>}, permitiendo especificar un rango mínimo y máximo de procesos permitidos. El gestor iniciará tantos procesos como pueda dentro de ese rango sin abortar si no se alcanza el máximo \cite{open_mpi_2026}.}: permite proporcionar al sistema un conjunto o rango de procesos permitidos, indicando a MPI que inicie el mayor número de procesos posible basándose en la disponibilidad real del clúster sin generar un error fatal \cite{mpi_forum_2025}. Esta estrategia es particularmente útil en entornos compartidos o con colas de trabajos dinámicas donde las ranuras pueden estar ocupadas parcialmente.

Un error lógico muy común entre los desarrolladores principiantes ocurre al interactuar con la salida estándar (\texttt{stdout}) inmediatamente después de incrementar el valor de \texttt{-np}. Es erróneo asumir que las sentencias de impresión aparecerán en la pantalla ordenadas según el rango de cada proceso, pues todos los procesos iniciados compiten simultáneamente por la consola, generando una salida impredecible o intercalada \cite{pacheco_2011}. Para evitarlo, la impresión siempre debe ser gestionada explícitamente desde el código, delegándola a un único proceso (comúnmente el rango 0). En ese sentido, la organización correcta del número de procesos no solo evita errores de lanzamiento, sino que también ayuda a mantener un comportamiento estable y predecible durante la ejecución paralela.
